\cotefile{162-5}
\scanpage{4}
\mspage{\handwritten{2}}

On a aussi un foncteur évident
\[ T:(\mathrm{Ssimpl})\longrightarrow(\mathrm{Cat}), \]
en associant à tout ensemble semi-simplicial $X$ la catégorie $T(X)=\handwritten{\Delta}_{/X}$
des simplexes sur $X$, dont l'ensemble des objets est la réunion d\unsure{\deleted{c}i|\deleted{i}c}sjointe
des $X_n$ .... (c'est une catégorie fibrée sur la catégorie $\handwritten{\Delta}$ des simplexes
types, à fibres les catégories discrètes définies par les $X_n$). Ceci posé
Quillen prouve que pour tout $X$, $ST(X)$ est isomorphe canoniquement à $X$ da
la catégorie homotopique construite avec (Ssimpl), et que pour toute $C$,
la catégorie $TS(C)$ est \begin{tabular}[b]{@{}l@{}}\inserted{canoniquement}\\[-.55ex]\begin{tikzpicture}[x=1em,y=1em,baseline=0pt]\draw[line width=.45pt] (0,-.7) -- (0,-.35);\draw[line width=.45pt] (.03,-.23) .. controls (.1,.2) and (.35,.38) .. (.6,.4);\draw[line width=.45pt] (.7,.41)--(1.4,.41);\draw[line width=.45pt] (1.5,.4) .. controls (3,.08) and (4,.45) .. (5.4,.15) .. controls (6.4,-.05) and (7.6,-.2) .. (8.2,.1) -- (8.2,.95);\end{tikzpicture}\end{tabular} ``homotopiquement équivalente à $C$'' \deleted{dans le sens} i.e
\inserted{\unsure{c|è}a\unsure{n|a}oniquement} isomorphe à $C$ dans la catégorie quotient de (Cat) obtenue en inversant le
foncteurs qui sont des homotopismes. \deleted{Il en résulte} Ces isomorphismes son
fonctoriels en $X$. I\rlap{\begin{tikzpicture}[x=1em,y=1em,baseline=0pt]\draw[line width=.045em] (.02,.95)--(.36,.95);\draw[line width=.045em] (.18,.02)--(.18,-.19);\end{tikzpicture}}l en résulte formellement qu'un morphisme $f:X\longrightarrow Y$ dans
(Ssimpl) est un homotopisme si et seulement si il en est ainsi de $T(f):T($
$\longrightarrow T(Y)$, d'où des foncteurs $S':(\mathrm{Cat})'\longrightarrow(\mathrm{Ssimpl})'$ et $T':(\mathrm{Ssimpl})'\longrightarrow(\mathrm{C}$
entre les catégories ``homotopiques'', \begin{tabular}[b]{@{}l@{}}\inserted{construites avec (Cat) resp.(Ssimpl),}\\[-.7ex]\begin{tikzpicture}[x=1em,y=1em,baseline=0pt]\draw[line width=.45pt] (.5,.5) .. controls (-.05,.4) and (-.15,-.1) .. (.1,-.7);\draw[line width=.45pt] (16,.15) .. controls (16.25,.4) and (16.45,.62) .. (16.5,.9);\end{tikzpicture}\end{tabular} qui sont quasi-inverses l'un de l'au
re.

\deleted{d}De plus, Quillen construit un isomorphisme canonique et fonctoriel
dans (Cat)' entre $C$ et la catégorie opposée $C^{o}$, ou ce qui revient au
même, un isomorphisme canonique et fonctoriel dans (Ssimpl)' entre $S(C)$ e
$S(C^{o})$. \deleted{Cela lui permet alors, pour} La définition est telle que le
foncteur induit sur les systèmes locaux sur $C$ trahs\unsure{\deleted{o}f|\deleted{f}o}\unsure{\deleted{f}o|\deleted{o}f}rme le foncteur
contravariant $F$ sur $C$, transformant toute flêche en flêche inversible, en
le foncteur covariant (i.e. contravariant sur $C^{o}$) ayant mêmes valeurs sur
les objets, et obtenu sur les flêches en remplaçant $F(u)$ par $F(u)^{-1}$; en
d'autres termes, l'effet de l'homotopisme de Quillen sur les group\unsure{\deleted{i}o|\deleted{o}i}ides
fondamentaux est l'isomorphisme évident entre les groupoides fondamentaux
de $C$ et de $C^{o}$, \deleted{qui est l'id \unsure{en|em}} compte tenu que le deuxième est l'opposé d\unsure{u}

\margin{\handwritten{3}}
